\documentclass{article}
\usepackage{graphicx} % Required for inserting images
\usepackage{amsmath}
\usepackage{tabularx}
\usepackage{pdflscape}
\usepackage{booktabs}
\usepackage[margin=1in]{geometry}
\usepackage{tabularray}
\usepackage{listings}
\usepackage{xcolor}
\usepackage{float}
\usepackage{minted}


% Definizione dei colori
\definecolor{codegreen}{rgb}{0,0.6,0}
\definecolor{codegray}{rgb}{0.5,0.5,0.5}
\definecolor{codepurple}{rgb}{0.58,0,0.82}
\definecolor{backcolour}{rgb}{0.95,0.95,0.92}
\definecolor{bluekeywords}{rgb}{0.13,0.13,1}

% Definizione dello stile SystemVerilog
\lstdefinestyle{mystyle}{
    backgroundcolor=\color{backcolour},   
    commentstyle=\color{codegreen},
    keywordstyle=\color{bluekeywords}\bfseries,
    numberstyle=\tiny\color{codegray},
    stringstyle=\color{codepurple},
    basicstyle=\ttfamily\footnotesize,
    breakatwhitespace=false,         
    breaklines=true,                 
    captionpos=b,                    
    keepspaces=true,                 
    numbers=left,                    
    numbersep=5pt,                  
    showspaces=false,                
    showstringspaces=false,
    showtabs=false,                  
    tabsize=2,
    language=Verilog,
    % Aggiunta di keyword specifiche per SystemVerilog presenti nel tuo codice
    morekeywords={logic, string, task, endtask, assert, enum, struct, typedef, byte, int, longint, shortint, bit, logic, reg, input, output, inout, initial, begin, end, module, endmodule, case, endcase, default, if, else, for, timescale, finish, display}
}

\lstset{style=mystyle}

\title{Processor Verification - Lab 4 \\ Definition Specification Document}
\author{Lorenzo Pattaro Zonta}

\begin{document}

\maketitle

\section{Fetch Unit}

\subsection{Functional Description}
The fetch unit maintains a program counter (PC) and outputs the instruction located at the current PC address. It supports normal sequential execution and conditional branching via an alternate branch address.

\subsection{Module Interface}

\begin{table}[H]
\centering
    \begin{tabular}{|l|l|l|p{8cm}|}
    \hline
    \textbf{Signal} & \textbf{Direction} & \textbf{Width} & \textbf{Description} \\
    \hline
    clk         & input  & 1   & Clock signal. \\
    rst\_n      & input  & 1   & Active-low synchronous reset. \\
    pc\_en      & input  & 1   & Enables PC increment (sequential fetch). \\
    branch\_en  & input  & 1   & Enables branch to branch\_addr. \\
    branch\_addr& input  & 8   & Branch target address. \\
    instr       & output & 16  & Instruction word fetched from memory. \\
    \hline
    \end{tabular}
\end{table}

\subsection{Behavioral Description}
An instruction memory is initialized with dummy values. On each clock edge, the PC updates to either the branch address (if branch is enabled), the next sequential address (PC + 2 if \texttt{pc\_en} is asserted), or holds its value. The output instruction is read from memory at the updated PC. On reset, the PC is cleared to zero and the instruction at address zero is output.


\subsection{List of assertions}
This is the list of the assertions for the fetch unit: 
\begin{itemize}
    \item Reset: after a reset cycle, PC must be zero
    \begin{minted}[fontsize=\small,breaklines]{systemverilog}
        property p_reset_pc_zero;
          @(posedge clk) !rst_n |=> (dut.pc == 8'h00);
        endproperty
        assert property (p_reset_pc_zero)
          else $error("PC not zero after reset");
    \end{minted}
    \item \begin{minted}[fontsize=\small,breaklines]{systemverilog}
      property p_branch_priority_load;
    @(posedge clk) disable iff (!rst_n)
      $past(branch_en) |-> (dut.pc == $past(branch_addr));
  endproperty

  assert property (p_branch_priority_load)
    else $error("Branch priority load failed");
    \end{minted}
    \item \begin{minted}[fontsize=\small,breaklines]{systemverilog}
      property p_branch_priority_not_increment;
    @(posedge clk) disable iff (!rst_n)
      $past(branch_en) |-> (dut.pc != ($past(dut.pc) + 8'd2));
  endproperty

  assert property (p_branch_priority_not_increment)
    else $error("Branch priority not-increment failed");
    \end{minted}
    \item \begin{minted}[fontsize=\small,breaklines]{systemverilog}
      property p_pc_increment;
    @(posedge clk) disable iff (!rst_n)
      ($past(pc_en) && !$past(branch_en)) |-> (dut.pc == ($past(dut.pc) + 8'd2));
  endproperty
  assert property (p_pc_increment)
    else $error("PC increment failed");
    \end{minted}
    \item \begin{minted}[fontsize=\small,breaklines]{systemverilog}
      property p_pc_hold;
    @(posedge clk) disable iff (!rst_n)
      (!$past(pc_en) && !$past(branch_en)) |-> (dut.pc == $past(dut.pc));
  endproperty

  assert property (p_pc_hold)
    else $error("PC hold failed");
    \end{minted}
    \item \begin{minted}[fontsize=\small,breaklines]{systemverilog}
      property p_instr_matches_mem;
    @(posedge clk) disable iff (!rst_n)
      (instr == dut.mem[dut.pc]);
  endproperty

  assert property (p_instr_matches_mem)
    else $error("Instruction does not match memory at PC");
    \end{minted}
    \item \begin{minted}[fontsize=\small,breaklines]{systemverilog}
      property p_pc_in_range;
    @(posedge clk) disable iff (!rst_n)
      (dut.pc inside {[8'h00:8'hFF]});
  endproperty
  assert property (p_pc_in_range)
    else $error("PC out of range");
    \end{minted}
\end{itemize}

\section{Decode Unit}

\subsection{Functional Description}
The decode unit receives a 16-bit instruction and splits it into opcode, destination register, source register, and immediate fields. It models a fixed two-cycle decode latency and performs a simple data hazard check that can stall decoding.

\subsection{Module Interface}
\begin{table}[H]
\centering
    \begin{tabular}{|l|l|l|p{8cm}|}
    \hline
    \textbf{Signal} & \textbf{Direction} & \textbf{Width} & \textbf{Description} \\
    \hline
    clk          & input  & 1   & Clock signal. \\
    rst\_n       & input  & 1   & Active-low synchronous reset. \\
    instr\_valid & input  & 1   & Indicates the input instruction is valid. \\
    instr        & input  & 16  & Instruction word from fetch stage. \\
    decode\_done & output & 1   & Pulses when decode latency completes. \\
    opcode       & output & 4   & Opcode field (bits 15:12). \\
    rd           & output & 4   & Destination register (bits 11:8). \\
    rs           & output & 4   & Source register (bits 7:4). \\
    imm          & output & 4   & Immediate field (bits 3:0). \\
    hazard\_stall& output & 1   & High when a data hazard is detected. \\
    \hline
    \end{tabular}
\end{table}

\subsection{Behavioral Description}
Instruction acceptance occurs when \texttt{instr\_valid}=1 and \texttt{hazard\_stall}=0 in the next cycle. On acceptance, the fields \texttt{opcode, rd, rs,} and \texttt{imm} capture instr[15:12], instr[11:8], instr[7:4], and instr[3:0], respectively. After an instruction is accepted in cycle $N$, \texttt{decode\_done} must assert exactly in cycle $N+2$ and it must be a single-cycle pulse.

A hazard is detected when the destination register \texttt{rd} of the previously accepted instruction matches the source register \texttt{rs} of the current instruction. When this hazard occurs, \texttt{hazard\_stall asserts}, the current instruction is not accepted until the stall clears, and \texttt{decode\_done} must not assert during the stall window. While \texttt{instr\_valid}=1 and \texttt{hazard\_stall}=1, the decoded fields (\texttt{opcode, rd, rs, imm}) must retain their previous values and only update upon acceptance.

For back-to-back valid instructions with no hazard, each instruction is accepted immediately in its cycle, and each acceptance produces a \texttt{decode\_done} pulse two cycles later, which may result in consecutive \texttt{decode\_done} pulses.

When \texttt{rst\_n}=0, all outputs are cleared to zero, and no instruction is accepted or completed until reset is released.

% \section{Objectives of the verification}
% The primary objectives of this verification are:
% \begin{itemize}
%     \item 
% \end{itemize}


% \section{Verification Plan Details}

% \begin{center}
% \begin{longtblr}[
%   label = {tab:testo},
%   caption = {Verification Plan Metrics}
% ]{
%   colspec = {c c c }, 
%   rowhead = 1, 
%   hlines,      
%   vlines,      
%   row{1} = {font=\bfseries}, 
% }  
%     Metric & Description & Goal \\
%      & & \\
% \end{longtblr}
% \end{center}


% \section{Results and Analysis}


% \newpage
% \begin{center}
%     \begin{longtblr}[
%   label = {tab:test_table},
%   caption = {Results of testbench simulations}
% ]{
%   colspec = {c c c c}, % 'X' funziona come in tabularx
%   rowhead = 1, % Ripete la prima riga (intestazione) su ogni pagina
%   hlines,      % Disegna tutte le linee orizzontali automaticamente
%   vlines,      % Disegna tutte le linee verticali
%   row{1} = {font=\bfseries}, % Rende l'intestazione grassetto
% }  
%     Design & Test Pass Rate & Total \# of errors & Requirement(s) Failed  \\
%     & & & \\
% \end{longtblr}
% \end{center}

% \section{Tools used}
% The verification process has been performed using the following tools:
% \begin{itemize}
%     \item QuestaSim
% \end{itemize}

\section{Conclusion}



\newpage
\section{Appendix}
\subsection{Number of hours required}
For the completion of this laboratory, it has been needed x hours. 
\subsection{Appendix A}
% \lstinputlisting[language=Bash, firstnumber=1]{script_questasim.sh}

\end{document}
